\SourceLeaf{443}%
Zusammentreffen zweier Linien gebildet werden. Offenbar bleiben daher auch Linien, die auf der einen Fläche kürzeste Linien waren, bei der Übertragung kürzeste Linien. Hieraus folgt leicht, dass, wenn einem beliebigen durch kürzeste Linien gebildeten Polygon auf der ersten Fläche die Figur der Zenite auf der Hilfskugel mit dem Flächeninhalt $=A$ entspricht und demselben Polygon bei seiner Übertragung durch Abwicklung auf die zweite Fläche eine Figur der Zenite auf der Hilfskugel mit dem Flächeninhalt $=A'$ entspricht, stets
\[
A=A'
\]
gilt. Obwohl dieser Beweis ursprünglich voraussetzt, dass die Figuren durch kürzeste Linien begrenzt sind, erkennt man leicht, dass er unabhängig von der Begrenzung allgemein gültig ist. Denn wenn der Lehrsatz von der Anzahl der Seiten unabhängig ist, kann man sich für jedes Polygon, dessen Seiten alle oder zum Teil keine kürzesten Linien sind, ein anderes mit unendlich vielen Seiten denken, die alle kürzeste Linien sind.

Ferner ist klar, dass jede Figur bei der Übertragung durch Abwicklung auch ihren Flächeninhalt behält.

Es ist nun hier von $4$ Figuren die Rede:
\begin{enumerate}
\item[1)] eine beliebige Figur auf der ersten Fläche,
\item[2)] die Figur auf der Hilfskugel, die den Zeniten der ersten Figur entspricht,
\item[3)] die Figur auf der zweiten Fläche, die durch Abwicklung von Nr. $1$ entsteht,
\item[4)] die Figur auf der Hilfskugel, die den Zeniten von Nr. $3$ entspricht.
\end{enumerate}
\noindent\hspace*{1.2em}Nach dem, was wir bewiesen haben, haben also $2$ und $4$ denselben Flächeninhalt, ebenso wie $1$ und $3$. Nimmt man nun diese Figuren als unendlich klein an, so ist der Quotient aus $2$ und $1$ das Krümmungsmaß an der Stelle der ersten gekrümmten Fläche und ebenso der Quotient aus $4$ und $3$ das Krümmungsmaß an der Stelle der zweiten. Daraus folgt der wichtige Lehrsatz, dass
\begin{quote}
»bei der Übertragung der Flächen durch Abwicklung das Krümmungsmaß an jeder Stelle unverändert bleibt«.
\end{quote}
\noindent Dasselbe gilt also für das Produkt des größten und des kleinsten Krümmungsradius.

Für die Ebene ist das Krümmungsmaß offenbar überall $=0$. Hieraus
